% GENERATED FILE. Edit canonical lessons, then run npm run generate:books.
% canonical-source-hash: fnv1a64:b3433c9792c981e3
% canonical-lessons: BN-C264-gelam, BN-C264-chhilam, BN-C264-ghota, BN-W264-letter-gha, BN-R264-what-happened

\chapter{What Happened, and Where I Was}
\label{ch:bn-what-happened-and-where-i-was}
\hlchaptermodality{264}

\begin{chapteropening}
\textbf{By the end of this chapter:} \emph{I can say what happened with the past by person (gelām, gelô, gelen), say where I was (chhilām), and read and write \bn{ঘ} in \bn{ঘটা}.}

Complete BN-R264-what-happened: say what happened and where you were, in the past.
\end{chapteropening}

\section[\texorpdfstring{\bn{আমি} \bn{গেলাম} (āmi gelām)}{আমি গেলাম (āmi gelām)}]{\bn{আমি} \bn{গেলাম} (āmi gelām) — the past, person by person}

\label{lesson:BN-C264-gelam}



\begin{quote}
Before the new one: say the Bengali for a female singer and a female dancer. (\emph{gāyikā}, \emph{nôrtôkī}.) Then the line from your first meeting: \textbf{\bn{আলাপ} \bn{করে} \bn{ভালো} \bn{লাগলো}}, \textquotedblleft{}pleased to meet you\textquotedblright{}.

Its last word is a past, \textquotedblleft{}it felt good\textquotedblright{}. This book now writes that ending \textbf{-\bn{ল}}: \textbf{\bn{লাগল}} (\emph{lāglô}). You will meet \textbf{-\bn{লো}} in print too; they sound the same.
\end{quote}

\subsection*{You'll want to know: the past, person by person}
\begin{quote}
\textbf{\bn{কাল} \bn{আমি} \bn{বাজারে} \bn{গেলাম।}}
\end{quote}

\begin{quote}
\emph{kāl āmi bājāre gelām.} — \textquotedblleft{}Yesterday I went to the market.\textquotedblright{}
\end{quote}

With a past, \textbf{\bn{কাল}} is yesterday. \textbf{\bn{গেলাম}} is \textbf{\bn{যাওয়া}} in the past: the stem \textbf{\bn{গে}-}, then \textbf{-\bn{ল}-}, then the person.

\begin{itemize}
\raggedright
  \item \textbf{\bn{আমি}} / \textbf{\bn{আমরা}} \textbf{\bn{গেলাম}} (\emph{gelām}) — I went, we went
  \item \textbf{\bn{তুমি}} / \textbf{\bn{তোমরা}} \textbf{\bn{গেলে}} (\emph{gele}) — you went, to a friend or to several
  \item \textbf{\bn{সে}} / \textbf{\bn{তারা}} \textbf{\bn{গেল}} (\emph{gelô}) — he or she went; they went
  \item \textbf{\bn{আপনি}} / \textbf{\bn{তিনি}} \textbf{\bn{গেলেন}} (\emph{gelen}) — you went, he or she went, with respect
\end{itemize}

They are the present's four slots: \textbf{\bn{যাই}}, \textbf{\bn{যাও}}, \textbf{\bn{যান}}, and one not written before, \textbf{\bn{সে} \bn{যায়}} (\emph{she jāy}), \textquotedblleft{}he or she goes\textquotedblright{}: \textbf{-\bn{য়}} after a vowel, \textbf{-\bn{এ}} after a consonant, \textbf{\bn{সে} \bn{বসে}}.

\subsection*{You'll want to know: the past stems}
Most verbs put \textbf{-\bn{ল}-} on the stem: \textbf{\bn{বসা}} $\to$ \textbf{\bn{বসলাম}} (\emph{bôslām}), \textbf{\bn{দেখা}} $\to$ \textbf{\bn{দেখলাম}} (\emph{dekhlām}), with a plain \emph{e} for the \emph{æ} of \textbf{\bn{দেখা}}. An \textbf{\bn{এ}} said \emph{e} rises to \textbf{\bn{ি}}, an \textbf{\bn{ও}} to \textbf{\bn{ু}}: \textbf{\bn{কেনা}} $\to$ \textbf{\bn{কিনলাম}} (\emph{kinlām}), \textbf{\bn{ফেরা}} $\to$ \textbf{\bn{ফিরলাম}} (\emph{phirlām}), \textbf{\bn{শোনা}} $\to$ \textbf{\bn{শুনলাম}} (\emph{shunlām}).

Eight verbs change the stem:

\begin{itemize}
\raggedright
  \item \textbf{\bn{যাওয়া}} $\to$ \textbf{\bn{গেলাম}} (\emph{gelām}) — went
  \item \textbf{\bn{আসা}} $\to$ \textbf{\bn{এলাম}} (\emph{elām}), \textbf{\bn{সে} \bn{এল}} (\emph{she elô}) — came
  \item \textbf{\bn{খাওয়া}} $\to$ \textbf{\bn{খেলাম}} (\emph{khelām}) — ate, drank
  \item \textbf{\bn{নেওয়া}} $\to$ \textbf{\bn{নিলাম}} (\emph{nilām}) — took
  \item \textbf{\bn{দেওয়া}} (of \textbf{\bn{জল} \bn{দেওয়া}}) $\to$ \textbf{\bn{দিলাম}} (\emph{dilām}) — gave
  \item \textbf{\bn{হওয়া}} $\to$ \textbf{\bn{হলাম}} (\emph{hôlām}); \textbf{\bn{কী} \bn{হল}?} (\emph{kī hôlô?}) — what happened?
  \item \textbf{\bn{চাওয়া}} $\to$ \textbf{\bn{চাইলাম}} (\emph{chāilām}), \textbf{\bn{সে} \bn{চাইল}} (\emph{she chāilô}) — wanted, asked for
  \item \textbf{\bn{পাওয়া}} $\to$ \textbf{\bn{পেলাম}} (\emph{pelām}) — got, found
\end{itemize}

\subsection*{Your turn}
Say these aloud:

\begin{itemize}
\raggedright
  \item \textquotedblleft{}Yesterday I went to the market\textquotedblright{} — \emph{kāl āmi bājāre gelām}
  \item \textquotedblleft{}In the evening she came back home\textquotedblright{} — \emph{sôndhyāy she bāṛi phirlô}
\end{itemize}

\begin{tcolorbox}[breakable,colback=teal!4,colframe=teal!35!black,title={Before you move on}]
\textquotedblleft{}I went\textquotedblright{}? (\textbf{\bn{গেলাম}}.) \textquotedblleft{}He came\textquotedblright{}? (\textbf{\bn{এল}}.) \textquotedblleft{}You went\textquotedblright{}, with respect? (\textbf{\bn{গেলেন}}.) \textquotedblleft{}What happened?\textquotedblright{} (\textbf{\bn{কী} \bn{হল}?})
\end{tcolorbox}
\section[\texorpdfstring{\bn{আমি} \bn{ছিলাম} (āmi chhilām)}{আমি ছিলাম (āmi chhilām)}]{\bn{আমি} \bn{ছিলাম} (āmi chhilām) — was, and was not}

\label{lesson:BN-C264-chhilam}



\begin{quote}
Say \textquotedblleft{}I went\textquotedblright{} and \textquotedblleft{}she came\textquotedblright{}. (\emph{gelām}, \emph{she elô}.) Then \textbf{\bn{আছি}}, \textquotedblleft{}I am, I am there\textquotedblright{}: this lesson puts it in the past.
\end{quote}

\subsection*{You'll want to know: \bn{ছিলাম}}
\begin{quote}
\textbf{\bn{কাল} \bn{আমি} \bn{বাড়িতে} \bn{ছিলাম।}}
\end{quote}

\begin{quote}
\emph{kāl āmi bāṛite chhilām.} — \textquotedblleft{}Yesterday I was at home.\textquotedblright{}
\end{quote}

\textbf{\bn{আছি}} has a past of its own, \textbf{\bn{ছিলাম}}, with the endings of the last lesson:

\begin{itemize}
\raggedright
  \item \textbf{\bn{আমি}} / \textbf{\bn{আমরা}} \textbf{\bn{ছিলাম}} (\emph{chhilām}) — I was, we were
  \item \textbf{\bn{তুমি}} \textbf{\bn{ছিলে}} (\emph{chhile}) — you were
  \item \textbf{\bn{সে}} / \textbf{\bn{তারা}} \textbf{\bn{ছিল}} (\emph{chhilô}) — he or she was; they were
  \item \textbf{\bn{আপনি}} / \textbf{\bn{তিনি}} \textbf{\bn{ছিলেন}} (\emph{chhilen}) — with respect
\end{itemize}

In the present, Bengali says \textquotedblleft{}is\textquotedblright{} with no verb at all: \textbf{\bn{সে} \bn{ভালো}}. In the past it needs one: \textbf{\bn{সে} \bn{ক্লান্ত} \bn{ছিল}} (\emph{she klāntô chhilô}), \textquotedblleft{}she was tired\textquotedblright{}. \textbf{\bn{ছিল}} also says \textquotedblleft{}there was\textquotedblright{}: \textbf{\bn{গ্রামে} \bn{একটা} \bn{বড়} \bn{পুকুর} \bn{ছিল}} (\emph{grāme ekṭā bôṛo pukur chhilô}), \textquotedblleft{}there was a big pond in the village\textquotedblright{}.

\subsection*{You'll want to know: \bn{ছিল} \bn{না}}
\textbf{\bn{নেই}}, \textquotedblleft{}there is not\textquotedblright{}, which you know, has no past of its own. Its past is \textbf{\bn{ছিল} \bn{না}}, with \textbf{\bn{না}} after the verb, as in the present:

\begin{quote}
\textbf{\bn{আজ} \bn{বাজারে} \bn{মাছ} \bn{নেই।} \bn{কালও} \bn{ছিল} \bn{না।}}
\end{quote}

\begin{quote}
\emph{āj bājāre māchh nei. kālo chhilô nā.} — \textquotedblleft{}There is no fish at the market today. There was none yesterday either.\textquotedblright{}
\end{quote}

\subsection*{Your turn}
Say these aloud:

\begin{itemize}
\raggedright
  \item \textquotedblleft{}Yesterday I was at home\textquotedblright{} — \emph{kāl āmi bāṛite chhilām}
  \item \textquotedblleft{}She was tired\textquotedblright{} — \emph{she klāntô chhilô}
  \item \textquotedblleft{}There was a pond in the village, but no river\textquotedblright{} — \emph{grāme ekṭā pukur chhilô, kintu nôdī chhilô nā}
\end{itemize}

\begin{tcolorbox}[breakable,colback=teal!4,colframe=teal!35!black,title={Before you move on}]
\textquotedblleft{}I was\textquotedblright{}? (\textbf{\bn{ছিলাম}}.) \textquotedblleft{}They were\textquotedblright{}? (\textbf{\bn{ছিল}}.) \textquotedblleft{}There was not\textquotedblright{}? (\textbf{\bn{ছিল} \bn{না}}.) What does the past need that the present leaves out? (The verb: \textbf{\bn{সে} \bn{ভালো}}, but \textbf{\bn{সে} \bn{ক্লান্ত} \bn{ছিল}}.)
\end{tcolorbox}
\section[\texorpdfstring{\bn{ঘটা} (ghôṭā)}{ঘটা (ghôṭā)}]{\bn{ঘটা} (ghôṭā) — to happen}

\label{lesson:BN-C264-ghota}



\begin{quote}
Before the new one: ask \textquotedblleft{}What happened?\textquotedblright{} the everyday way. (\textbf{\bn{কী} \bn{হল}?}, \emph{kī hôlô?})
\end{quote}

\subsection*{You'll want to know: \bn{ঘটা}}
\textbf{\bn{ঘটা}} — \emph{ghôṭā} — \textquotedblleft{}to happen\textquotedblright{}.

\begin{quote}
\textbf{\bn{কাল} \bn{রাতে} \bn{পাড়ায়} \bn{কী} \bn{ঘটল}?}
\end{quote}

\begin{quote}
\emph{kāl rāte pāṛāy kī ghôṭlô?} — \textquotedblleft{}What happened in the neighbourhood last night?\textquotedblright{}
\end{quote}

\textbf{\bn{কী} \bn{হল}?} is the everyday question. \textbf{\bn{ঘটা}} is for something that takes place, an event, and it suits a story. Its past is regular: \textbf{\bn{ঘটল}} (\emph{ghôṭlô}), \textquotedblleft{}it happened\textquotedblright{}.

\begin{quote}
\textbf{\bn{এটা} \bn{কখন} \bn{ঘটল}?}
\end{quote}

\begin{quote}
\emph{eṭā kôkhôn ghôṭlô?} — \textquotedblleft{}When did this happen?\textquotedblright{}
\end{quote}

\begin{cousinweb}
\textbf{\bn{ঘটা}} (\emph{ghôṭā}) is usually traced to the Sanskrit verb \emph{ghaṭate}, which CDIAL (4407) glosses \textquotedblleft{}is busy with\textquotedblright{} and, in a later text, \textquotedblleft{}happens\textquotedblright{}. CDIAL lists no Bengali form under that entry, so the line from one to the other is the usual tracing, not a recorded row.
\end{cousinweb}

\subsection*{Your turn}
Say these aloud:

\begin{itemize}
\raggedright
  \item \emph{ghôṭā}, then its past, \emph{ghôṭlô}
  \item \textquotedblleft{}What happened in the neighbourhood last night?\textquotedblright{} — \emph{kāl rāte pāṛāy kī ghôṭlô?}
  \item \textquotedblleft{}When did this happen?\textquotedblright{} — \emph{eṭā kôkhôn ghôṭlô?}
\end{itemize}

\begin{tcolorbox}[breakable,colback=teal!4,colframe=teal!35!black,title={Before you move on}]
What does \textbf{\bn{ঘটা}} mean? (\textquotedblleft{}To happen\textquotedblright{}.) Its past? (\textbf{\bn{ঘটল}}, \emph{ghôṭlô}.) Which question is the everyday one? (\textbf{\bn{কী} \bn{হল}?})
\end{tcolorbox}
\section[\texorpdfstring{\bn{ঘ} (ghô)}{ঘ (ghô)}]{\bn{ঘ} (ghô) — the letter at the start of \bn{ঘটা}}

\label{lesson:BN-W264-letter-gha}



\begin{quote}
Before the new one: say \textbf{\bn{ঘটা}}, \textquotedblleft{}to happen\textquotedblright{}, and its past, \textbf{\bn{ঘটল}}. Then read \textbf{\bn{ব}} aloud (\emph{bô}); you will want its shape in a moment.
\end{quote}

\subsection*{Script you'll notice: \bn{ঘ}}
\textbf{\bn{ঘ}} — \emph{ghô}: a \emph{g} with a breath after it, as \textbf{\bn{খ}} (\emph{khô}) is a \emph{k} with a breath. It opens \textbf{\bn{ঘটা}}.

In the print of this book, \textbf{\bn{ঘ}} has the frame of \textbf{\bn{ব}}: a headline along the top, an upright stem on the right, the full height of the letter, a pointed tip at the lower left, and a curve that joins that tip to the foot of the stem. Where \textbf{\bn{ব}} has one straight diagonal, \textbf{\bn{ঘ}} has two pieces. Under the left end of the headline hangs a rounded curl, open toward the stem, its end a short point aimed up and to the right. Below the curl, a short straight piece joins it to the pointed tip.

\subsection*{Writing: \bn{ঘ} — copy the printed shape}
Put your pen beside \textbf{\bn{ঘ}} and copy the shape you can see, slowly and large. Check the copy against the print: the headline, the stem, the curl with its point, the tip at the lower left, and the curve to the foot.

\begin{quote}
This book does not yet tell you where to start \bn{ঘ} or which way the pen goes. That is not written down here until it can be written down with a source. Copying the printed shape needs no such source, and it is how the shape becomes yours in the meantime.
\end{quote}

\subsection*{Your turn}
\begin{itemize}
\raggedright
  \item \emph{Find:} \bn{ঘ} at the front of \textbf{\bn{ঘটা}}
  \item \emph{Write it:} \bn{ব}, then \bn{ঘ}, from the print, naming each
  \item \emph{Write it:} \textbf{\bn{ঘটা}}, then \textbf{\bn{ঘটল}}
  \item \emph{Read it:} \textbf{\bn{ঘটল}} aloud: \emph{ghôṭlô}
\end{itemize}

\begin{tcolorbox}[breakable,colback=teal!4,colframe=teal!35!black,title={Before you move on}]
Which letter is this — \bn{ঘ}? (\emph{ghô}.) Which word did it come from? (\textbf{\bn{ঘটা}}.) How is it unlike \textbf{\bn{ব}}? (A curl hangs under its headline.)
\end{tcolorbox}
\section[(review)]{Review — what happened}

\label{lesson:BN-R264-what-happened}



\begin{quote}
Say \textquotedblleft{}Yesterday I went to the market\textquotedblright{}. (\emph{kāl āmi bājāre gelām}.) Then \textquotedblleft{}she came\textquotedblright{} and \textquotedblleft{}we ate\textquotedblright{}. (\emph{she elô}, \emph{āmrā khelām}.)
\end{quote}

\subsection*{Your turn}
\begin{itemize}
\raggedright
  \item \emph{Say it:} In the evening I came back home. I was tired.
  \item \emph{Say it:} What happened in the village yesterday?
  \item \emph{Say it:} Mira was at the station. Arun was not there.
  \item \emph{Say it:} There was a big tree in the neighbourhood.
  \item \emph{Say it:} We wanted tea, and we got tea.
  \item \emph{Write it:} \bn{ঘ}, then \textbf{\bn{ঘটল}}
\end{itemize}

\begin{tcolorbox}[breakable,colback=teal!4,colframe=teal!35!black,title={Before you move on}]
Which ending is I and we? (\textbf{-\bn{লাম}}: \textbf{\bn{গেলাম}}.) Which verbs gave \textbf{\bn{এল}} and \textbf{\bn{পেল}}? (\textbf{\bn{আসা}}, \textbf{\bn{পাওয়া}}.) \textquotedblleft{}There was not\textquotedblright{}? (\textbf{\bn{ছিল} \bn{না}}.) Which letter opens \textbf{\bn{ঘটা}}? (\textbf{\bn{ঘ}}.)
\end{tcolorbox}
